% =============================================================
%  CS102-F26 Digital Design and Computer Architecture I
%  Lecture 09 -- Timing of Combinational Circuits (Unit 9)
%  Follows Section 4.2 of A Student's Computer Architecture Book I.
%
%  Build the slides and the print handout:  npm run materials -- lecture-09
% =============================================================
\makeatletter
\def\input@path{{templates/latex/}{../}}
\makeatother

\documentclass[10pt]{wosm-slides}
\usepackage{cs102-lectures}[2026/10/19]

\title[Lecture 09 \textperiodcentered\ Timing]{Timing}
\subtitle{Lecture 09: propagation delay, contamination delay and glitches}

% ---------- Helpers for this lecture --------------------------
% Wider input spacing than the course default, so input labels fit.
\tikzset{gates/.append style={
  and gate/.append style={logic gate input sep=0.36cm},
  or gate/.append style={logic gate input sep=0.36cm}}}

% Timing diagrams, 1 ns per unit, from t = -1 to t = 7.
% \wave{<row>}{<label>}{<level in each ns, from t = -1 to 6>}
\newcommand*\wave[3]{%
  \pgfmathsetmacro\wy{-#1*0.52}
  \node[left, font=\scriptsize] at (-1.1,\wy+0.14) {#2};
  \foreach \v [count=\i from -1, remember=\v as \pv (initially X)] in {#3} {
    \draw[wosm-text, line width=0.9pt] (\i,\wy+\v*0.28) -- (\i+1,\wy+\v*0.28);
    \if\pv X\else
      \draw[wosm-text, line width=0.9pt] (\i,\wy+\pv*0.28) -- (\i,\wy+\v*0.28);
    \fi
  }}
% \timeaxis{<rows>}: tick marks under the last row.
\newcommand*\timeaxis[1]{%
  \pgfmathsetmacro\wy{-#1*0.52+0.22}
  \draw[wosm-muted, line width=0.4pt] (-1,\wy) -- (7,\wy);
  \foreach \t in {0,...,6}
    \draw[wosm-muted, line width=0.4pt] (\t,\wy) -- ++(0,-0.08)
      node[below, font=\tiny] {\t};
  \node[right, font=\tiny, text=wosm-muted] at (7,\wy) {ns};}
% \tmark{<time>}{<rows>}{<label>}: a dashed line across all rows.
\newcommand*\tmark[3]{%
  \draw[wosm-brick, dashed, line width=0.5pt] (#1,0.45) -- (#1,{-#2*0.52+0.22});
  \node[above, font=\scriptsize, text=wosm-brick] at (#1,0.4) {#3};}

% The example circuit f = xy' + x'yz + z' (book, Figure 4.5).
% \bookcircuit{<style of the longest path>}{<style of the shortest path>}
\newcommand*\bookcircuit[2]{%
  \begin{tikzpicture}[gates, scale=0.88, every node/.append style={transform shape}]
    \coordinate (left) at (0,0);
    \node[and gate, inputs=nn] (a1) at (3.6,2.55) {};
    \node[and gate, inputs=nn] (a2) at (3.6,0.95) {};
    \node[and gate, inputs=nn] (a3) at (5.3,0.45) {};
    \node[or gate, inputs=nn] (o1) at (7.0,1.6) {};
    \node[or gate, inputs=nn] (o2) at (8.7,0.95) {};
    \node[not gate] (ny) at ($(a1.input 2)+(-1.9,0)$) {};
    \node[not gate] (nx) at ($(a2.input 1)+(-1.9,0)$) {};
    \node[not gate] (nz) at (6.9,-0.7) {};
    \node (x) at (a1.input 1 -| left) {$x$};
    \node (y) at (a2.input 2 -| left) {$y$};
    \node (z) at (0,-0.7) {$z$};
    % x: straight to the first AND, and through a NOT to the second
    \draw[wire, #1] (x) -- (x -| 0.9,0);
    \draw[wire] (x -| 0.9,0) -- (a1.input 1);
    \draw[wire, #1] (x -| 0.9,0) |- (nx.input);
    \fill[wosm-text] (x -| 0.9,0) circle (1.5pt);
    \draw[wire, #1] (nx.output) -- (a2.input 1);
    % y: straight to the second AND, and through a NOT to the first
    \draw[wire] (y) -- (a2.input 2);
    \draw[wire] (y -| 0.5,0) |- (ny.input);
    \fill[wosm-text] (y -| 0.5,0) circle (1.5pt);
    \draw[wire] (ny.output) -- (a1.input 2);
    % z: to the third AND, and through a NOT to the last OR
    \draw[wire, #2] (z) -- (nz.input);
    \draw[wire] (4.6,-0.7) |- (a3.input 2);
    \fill[wosm-text] (4.6,-0.7) circle (1.5pt);
    \draw[wire, #1] (a2.output) -- ++(0.25,0) |- (a3.input 1);
    \draw[wire] (a1.output) -- ++(0.4,0) |- (o1.input 1);
    \draw[wire, #1] (a3.output) -- ++(0.25,0) |- (o1.input 2);
    \draw[wire, #1] (o1.output) -- ++(0.25,0) |- (o2.input 1);
    \draw[wire, #2] (nz.output) -- ++(0.4,0) |- (o2.input 2);
    \draw[wire] (o2.output) -- ++(0.35,0) node[right] {$f$};
    \node[edgelabel, above] at (a1.north) {$xy'$};
    \node[edgelabel, below=1pt] at (a2.south) {$x'y$};
    \node[edgelabel, below=1pt] at (a3.south) {$x'yz$};
  \end{tikzpicture}}
\tikzset{longest/.style={draw=wosm-brick, line width=1.4pt},
  shortest/.style={draw=wosm-accent, line width=1.4pt}}

% =============================================================
\begin{document}

\maketitle

\begin{frame}{Today}
  \tableofcontents
\end{frame}

% =============================================================
\section{Time in circuits}

\begin{frame}{On paper, there is no time}
  Take $f(x, y, z) = xy' + x'yz + z'$.

  \begin{itemize}
    \item Computing $f(1, 1, 0)$ by hand might take you 5 seconds. But you only give
      the answer when you are done, so whoever hears it knows it is right.
    \item A circuit cannot keep quiet until it is done. Its output wire always
      carries \emph{some} value, 0 or 1, even while its gates are still working.
  \end{itemize}

  \bigskip
  \begin{block}{The question for today}
    After the inputs change, how long until the output can be trusted? And what does
    it do in the meantime?
  \end{block}
\end{frame}

\begin{frame}{The example circuit}
  \begin{center}
    \bookcircuit{}{}
  \end{center}

  \begin{itemize}
    \item 0 and 1 are voltages, for example 0\,V and 5\,V. Keep the inputs at
      $x, y, z = 1, 1, 0$ and the output settles at $f = 1$.
    \item Then change the inputs to $0, 1, 1$. Here too $f(0, 1, 1) = 1$, so we expect
      the output to stay at 1. Does it?
  \end{itemize}
\end{frame}

\begin{frame}{Delays}
  \begin{itemize}
    \item A signal needs time to travel through wires and through the transistors
      inside each gate: a \alert{delay}. The longer the wire, the longer the delay.
    \item To keep things simple, as the book does: \alert{every gate takes exactly
      1\,ns}, and wires take no time.
    \item So a gate's output at time $t$ is what its inputs were at time
      $t - 1\,\mathrm{ns}$.
  \end{itemize}

  \bigskip
  \begin{remark}
    Real gates switch in far less than a nanosecond, and different gates have
    different delays. The method is the same: only the numbers change.
  \end{remark}
\end{frame}

% =============================================================
\section{Propagation and contamination delay}

\begin{frame}{Two delays}
  \begin{columns}[c, totalwidth=\textwidth]
    \column{0.48\textwidth}
      \begin{definition}[Propagation delay $t_{pd}$]
        The \emph{longest} time from a change at the inputs until the outputs have
        reached their final values.
      \end{definition}
      \begin{definition}[Contamination delay $t_{cd}$]
        The \emph{shortest} time from a change at the inputs until an output starts
        to change.
      \end{definition}
    \column{0.48\textwidth}
      \centering
      \begin{tikzpicture}[x=0.55cm]
        \draw[wosm-text, line width=0.9pt] (-1,0.32) -- (0,0.32) -- (0,0) -- (7,0);
        \node[left, font=\scriptsize] at (-1.1,0.15) {input};
        \draw[wosm-text, line width=0.9pt] (-1,-0.6) -- (2,-0.6);
        \draw[wosm-text, line width=0.9pt] (5,-0.6) -- (7,-0.6);
        \fill[wosm-rule] (2,-0.62) rectangle (5,-0.26);
        \draw[wosm-text, line width=0.9pt] (2,-0.6) -- (2,-0.28) -- (5,-0.28);
        \node[font=\tiny, text=wosm-muted] at (3.5,-0.44) {may change};
        \node[left, font=\scriptsize] at (-1.1,-0.45) {output};
        \draw[wosm-brick, dashed, line width=0.5pt] (0,0.6) -- (0,-0.9);
        \draw[wosm-accent, line width=0.5pt, {Stealth[length=3pt]}-{Stealth[length=3pt]}]
          (0,-1.0) -- node[below, font=\scriptsize] {$t_{cd}$} (2,-1.0);
        \draw[wosm-brick, line width=0.5pt, {Stealth[length=3pt]}-{Stealth[length=3pt]}]
          (0,-1.55) -- node[below, font=\scriptsize] {$t_{pd}$} (5,-1.55);
      \end{tikzpicture}
  \end{columns}

  \medskip
  Before $t_{cd}$ the output still shows the old answer; after $t_{pd}$ it shows the
  new one. In between, it may do anything.
\end{frame}

\begin{frame}{Finding the delays: paths}
  \begin{center}
    \bookcircuit{longest}{shortest}
  \end{center}

  \begin{itemize}
    \item \textcolor{wosm-brick}{Longest path}, the \alert{critical path}:
      $x \to$ NOT $\to$ AND $\to$ AND $\to$ OR $\to$ OR, 5 gates, so
      $t_{pd} = 5$\,ns.
    \item \textcolor{wosm-accent}{Shortest path}: $z \to$ NOT $\to$ OR, 2 gates, so
      $t_{cd} = 2$\,ns.
    \item The other four paths from an input to $f$ have 3 or 4 gates.
  \end{itemize}
\end{frame}

\begin{frame}{Gates with different delays}
  \begin{itemize}
    \item Each gate has its own $t_{pd}$ and $t_{cd}$, with $t_{cd} \le t_{pd}$.
    \item The circuit's $t_{pd}$ is the largest sum of gate $t_{pd}$'s along any path
      from an input to an output.
    \item Its $t_{cd}$ is the smallest sum of gate $t_{cd}$'s along any path.
  \end{itemize}

  \bigskip
  \begin{warning}
    With different gate delays, the critical path is not always the path with the
    most gates. Add up the delays of every path, or at least of every path that could
    be the longest.
  \end{warning}
\end{frame}

% =============================================================
\section{Glitches}

\begin{frame}{Watching the change, nanosecond by nanosecond}
  The inputs change from $x, y, z = 1, 1, 0$ to $0, 1, 1$ at $t = 0$.

  \medskip
  \begin{center}
    \begin{tabular}{@{}l|c|ccccccc@{}}
      $t$ (ns) & before & 0 & 1 & 2 & 3 & 4 & 5 & 6 \\
      \midrule
      $x'$ & 0 & 0 & 1 & 1 & 1 & 1 & 1 & 1 \\
      $z'$ & 1 & 1 & 0 & 0 & 0 & 0 & 0 & 0 \\
      $x'y$ & 0 & 0 & 0 & 1 & 1 & 1 & 1 & 1 \\
      $x'yz$ & 0 & 0 & 0 & 0 & 1 & 1 & 1 & 1 \\
      $xy' + x'yz$ & 0 & 0 & 0 & 0 & 0 & 1 & 1 & 1 \\
      $f$ & 1 & 1 & 1 & \alert{0} & \alert{0} & \alert{0} & 1 & 1 \\
    \end{tabular}
  \end{center}

  \medskip
  \begin{itemize}
    \item $z'$ falls after 1\,ns, so $f$ falls after 2\,ns: nothing else is 1 yet.
    \item The term $x'yz$ needs 4\,ns to rise, so $f$ is back at 1 only after 5\,ns.
  \end{itemize}
\end{frame}

\begin{frame}{The timing diagram}
  \begin{center}
    \begin{tikzpicture}[x=0.95cm]
      \wave{0}{$x$}{1,0,0,0,0,0,0,0}
      \wave{1}{$z$}{0,1,1,1,1,1,1,1}
      \wave{2}{$x'$}{0,0,1,1,1,1,1,1}
      \wave{3}{$z'$}{1,1,0,0,0,0,0,0}
      \wave{4}{$x'y$}{0,0,0,1,1,1,1,1}
      \wave{5}{$x'yz$}{0,0,0,0,1,1,1,1}
      \wave{6}{$xy' + x'yz$}{0,0,0,0,0,1,1,1}
      \wave{7}{$f$}{1,1,1,0,0,0,1,1}
      \timeaxis{8}
      \tmark{2}{8}{$t_{cd}$}
      \tmark{5}{8}{$t_{pd}$}
    \end{tikzpicture}
  \end{center}
  \wosmlabel{$y$ stays 1 throughout. Time 0 is when $x$ and $z$ change.}
\end{frame}

\begin{frame}{Glitches}
  \begin{definition}[Glitch]
    A \emph{glitch} is a short wrong value at an output while the circuit settles
    after its inputs change.
  \end{definition}

  \begin{itemize}
    \item Here $f$ should stay at 1, but it is 0 from 2\,ns to 5\,ns.
    \item Anyone who reads $f$ in that window gets the wrong answer.
    \item After $t_{pd}$ the output is always right. So the simplest rule is: wait at
      least $t_{pd}$ after the inputs change, then read the output.
  \end{itemize}

  \bigskip
  \wosmlabel{This needs $t_{pd}$ to be known, which is why we compute it.}
\end{frame}

\begin{frame}{A glitch from a single input}
  The multiplexer of Lecture 07, $y = x_1 s' + x_2 s$, with $x_1 = x_2 = 1$. Then
  $y = 1$ whatever $s$ is. Now switch $s$ from 1 to 0:

  \medskip
  \begin{columns}[c, totalwidth=\textwidth]
    \column{0.56\textwidth}
      \centering
      \begin{tikzpicture}[x=0.75cm]
        \wave{0}{$s$}{1,0,0,0,0,0,0,0}
        \wave{1}{$s'$}{0,0,1,1,1,1,1,1}
        \wave{2}{$x_2 s$}{1,1,0,0,0,0,0,0}
        \wave{3}{$x_1 s'$}{0,0,0,1,1,1,1,1}
        \wave{4}{$y$}{1,1,1,0,1,1,1,1}
        \timeaxis{5}
      \end{tikzpicture}
    \column{0.4\textwidth}
      \begin{itemize}
        \item $x_2 s$ turns off after 1\,ns.
        \item $x_1 s'$ waits for the NOT gate, and turns on only after 2\,ns.
        \item For 1\,ns neither door is open: $y$ glitches to 0.
      \end{itemize}
  \end{columns}
\end{frame}

\begin{frame}{Removing the glitch}
  \begin{columns}[c, totalwidth=\textwidth]
    \column{0.36\textwidth}
      \centering
      \begin{kmap}{4}{2}{x_1x_2}{s}{00,01,11,10}{0,1}
        \kmaprow{0}{0,0}\kmaprow{1}{0,1}\kmaprow{2}{1,1}\kmaprow{3}{1,0}
        \kgroup{2}{0}{3}{0}
        \kgroup[wosm-brick]{1}{1}{2}{1}
        \kgroup[wosm-text]{2}{0}{2}{1}
      \end{kmap}
    \column{0.6\textwidth}
      \begin{itemize}
        \item The change goes from cell 111 to cell 110: two 1-cells next to each other
          but in \alert{different groups}.
        \item Add the group that covers both, $x_1 x_2$: the consensus term we left
          out in Lecture 07.
        \item $y = x_1 s' + x_2 s + x_1 x_2$ is the same function, and $x_1 x_2$ stays
          1 while $s$ changes: no glitch.
      \end{itemize}
  \end{columns}

  \bigskip
  \begin{remark}
    The price is one more gate. Usually we do not pay it: we wait for $t_{pd}$ instead,
    and a clock tells every part of the computer when to read \unitref{11--14}.
  \end{remark}
\end{frame}

% =============================================================
\section{Why timing matters}

\begin{frame}{The speed limit}
  \begin{itemize}
    \item A result can only be used $t_{pd}$ after the inputs change. With
      $t_{pd} = 5$\,ns, that is at most one result every 5\,ns:
      \[ \frac{1}{5\,\mathrm{ns}} = 200 \text{ million results per second.} \]
    \item To make a circuit faster, shorten its critical path.
    \item In sequential circuits both delays matter: $t_{pd}$ sets how fast the clock
      can tick, and $t_{cd}$ how soon a stored value may be disturbed. These are the
      setup and hold time constraints \unitref{14}.
  \end{itemize}

  \bigskip
  \wosmlabel{The book: delays become more and more important as we build a real
  computer.}
\end{frame}

\begin{frame}{Try it}
  \begin{enumerate}
    \item In the example circuit, let NOT take 1\,ns, and AND and OR 2\,ns each.
      What are $t_{pd}$ and $t_{cd}$?\par
      \uncover<2->{\small $t_{pd} = 1 + 2 + 2 + 2 + 2 = 9$\,ns, still along
      $x \to$ NOT $\to$ AND $\to$ AND $\to$ OR $\to$ OR.
      $t_{cd} = 1 + 2 = 3$\,ns, along $z \to$ NOT $\to$ OR.}
    \item With 1\,ns gates, does the multiplexer glitch when $s$ goes from 0 to
      1?\par
      \uncover<3->{\small No. $x_2 s$ turns on after 1\,ns, and $x_1 s'$ turns off only
      after 2\,ns: for 1\,ns both are 1, so $y$ stays 1.}
    \item Why does a longer wire have a longer delay?\par
      \uncover<4->{\small The signal takes time to travel, and a longer wire takes
      longer to change its voltage.}
  \end{enumerate}
\end{frame}

% =============================================================
\section{Summary}

\begin{frame}{Summary}
  \begin{itemize}
    \item Real circuits take time: every gate and wire adds a delay, and the output
      always shows some value, right or not.
    \item Propagation delay $t_{pd}$: the longest path, the time until the output is
      final. Contamination delay $t_{cd}$: the shortest path, the time until it can
      first change.
    \item A glitch is a short wrong output while the circuit settles. Reading after
      $t_{pd}$ is always safe.
    \item A glitch from one input can be removed with the consensus term, at the cost
      of a gate.
    \item $t_{pd}$ limits how fast a circuit can be used.
  \end{itemize}
\end{frame}

\begin{frame}{Before the next lecture}
  \begin{itemize}
    \item Read Section~4.2 of the book, \emph{Timing}.
    \item For $g = xy + x'z$, list every path, and find $t_{pd}$ and $t_{cd}$ with 1\,ns
      gates. Which single input change can make $g$ glitch, and how do you remove it?
    \item In Logisim-evolution, build the example circuit as in the book: one input
      $\mathit{ctrl}$ drives $x$ directly and $z$ through a NOT gate, so both change
      together. Switch $\mathit{ctrl}$ from 1 to 0 and find the glitch in the timing
      diagram.
  \end{itemize}

  \begin{block}{Next lecture \unitref{10}}
    Sequential logic: finite state machines, Moore and Mealy machines.
  \end{block}
\end{frame}

\end{document}
